\subsection{Scaling law：把大模型赌博改造成可外推实验}

slides 17--23 从“多花计算通常能改善 loss”出发，说明 scaling law 的工程价值。旧流程在大模型上直接试大量超参数，失败试验本身就与最终训练一样昂贵；新流程在多个小规模点上拟合学习率、模型规模、数据量与性能之间的规律，再把 recipe 外推到目标规模。Chinchilla 的 IsoFLOP 分析进一步问：固定计算预算时，参数和 token 应怎样分配。Bitter Lesson 则提醒长期进步往往来自能随计算扩展的一般方法。

\lecturefigure{slides-images/slide-017.png}{官方 slide 17：计算量与性能之间的经验 scaling law}
\lecturefigure{slides-images/slide-018.png}{官方 slide 18：从大模型试错转向小规模 recipe 外推}
\lecturefigure{slides-images/slide-019.png}{官方 slide 19：用常数项与斜率比较架构的扩展表现}
\lecturefigure{slides-images/slide-020.png}{官方 slide 20：Chinchilla IsoFLOP 下的参数/数据最优分配}
\lecturefigure{slides-images/slide-021.png}{官方 slide 21：Bitter Lesson 与利用计算的一般方法}
\lecturefigure{slides-images/slide-022.png}{官方 slide 22：LLaMA 3 405B 训练 FLOPs、时间、成本与碳排估算}
\lecturefigure{slides-images/slide-023.png}{官方 slide 23：预训练阶段的数量级总结}

\begin{knowledgebox}{读图：斜率、常数与外推区间}
slide 19 中“Transformer 比 LSTM 更好”包含两个维度：在观察区间内起点更低，以及随 compute 增长下降得更快。只有少量点时，斜率估计很不稳定；跨越过大数量级外推还可能遇到数据质量、优化器、硬件和架构变化。slide 20 的 IsoFLOP 曲线也只回答固定训练计算下的 loss 最优，不直接回答部署延迟、显存、数据许可或推理成本最优。scaling law 是预算决策工具，不是取消机制研究的理由。
\end{knowledgebox}

\begin{importantbox}{Worked example：训练 FLOPs 怎样估算}
对 dense Transformer，课程采用近似式
\[
C_{\mathrm{train}}\approx 6ND,
\]
其中 $N$ 是参数量，$D$ 是训练 token 数，系数 6 粗略包含 forward 与 backward。以 $N=405\times10^9$、$D=15.6\times10^{12}$ 代入，得到约 $3.8\times10^{25}$ FLOPs。再除以 GPU 数、有效吞吐和时间单位，可估 GPU-hours。这里的关键不在最后一位数字，而在把“模型规模”转成可审计的资源假设。
\end{importantbox}

\begin{warningbox}{Bitter Lesson 不等于“只要堆卡”}
课程引用 Sutton 是为了反对无法随规模扩展的手工技巧，但 slide 18 同时强调 recipe、数据和小规模实验。若数据 mixture 错误、评估失真或系统利用率低，增加 compute 只会更昂贵地重复错误。\textbf{讲义提醒：}slide 22 的租赁单价、400 TFLOPS 平均吞吐和碳强度都是估算假设；它支持“前沿训练是系统级项目”，不证明任何机构的真实成本或环境影响精确等于表中数字。
\end{warningbox}

